\documentclass[11pt]{article}
\usepackage[letterpaper,margin=1in]{geometry}
\usepackage{booktabs}
\usepackage{longtable}
\usepackage{array}
\usepackage{xcolor}
\usepackage{microtype}
\setlength{\emergencystretch}{2em}
\usepackage[hidelinks]{hyperref}
\usepackage[capitalise,nameinlink]{cleveref}
\usepackage{caption}
\captionsetup{font=small,labelfont=bf}

% Protocol facts cite the code that ran: irregular-benchmarking commit 2593e29
% (src/, scripts/) and ABCG commit 29953fe (benchmarks/). Quotes were checked
% verbatim against the papers' texts on 2026-10-14. Numbers about the data come from
% data-10-14-2026/results/ (runs.csv, tables/) and the pilot output files.

\newcommand{\verdict}[1]{\textbf{#1}}
\newcommand{\pmrep}{\texttt{.pm} replay}

\title{Assessing the Data-Collection Protocol of the Five-Way MPI\\Halo-Exchange Campaign Against Published Benchmarking Guidance}
\author{Nicholas Bacon\\University of New Mexico}
\date{October 2026}

\begin{document}
\maketitle

\begin{abstract}
MPI performance measurements vary between separate launches of the same program, so
conclusions drawn from one launch, from unrandomized run orders, or from statistics that
assume normality can be wrong without any visible error. Published guidance from HPC and
from the broader systems and machine learning communities names these pitfalls and
states how to avoid them. However, applying that guidance to a large campaign involves
choices the guidance does not settle, such as the unit of replication and how to time an
exchange closed by a barrier. In this document we assess the data-collection protocol of
the October 2026 five-way halo-exchange campaign on LLNL's Dane, which ran 20,160 runs in
160 separate Slurm allocations, against 28 requirements drawn from Hoefler and Belli,
Hunold and Carpen-Amarie, related work, and the project's research plan. The protocol
meets 10, partly meets 13, and does not meet 5. It replicates every combination in 40
separate allocations per node count, randomizes run order within each allocation with a
recorded seed, compares strategies with medians and intervals resampled over
allocations, and keeps every outlier. It does not time single iterations, does not
interleave strategies within a run, does not test whether ranks differ before
summarizing them, fixed the number of allocations in advance instead of stopping on
interval width, and does not state why its six matrices were chosen. It also chose
iteration counts from pilot runs, some of them slowed by stalls, includes a barrier inside
every timed interval of the replay and SpMV benchmarks, and spans 23.7 hours rather than
weeks. We give
the assessment as a reusable checklist and list changes for the next campaign.
\end{abstract}

\section{Introduction}

The time an MPI program takes to communicate depends on factors the experimenter does
not control. Hunold and Carpen-Amarie found that separate calls to \texttt{mpirun} within
one job allocation produce means that ``are often not very large (3-5 percent), yet they
are statistically significant'', and concluded that ``a call to mpirun is a factor of the
experiment''~\cite{hunold2016reproducible}. Bhatele et al.\ attributed run-to-run
variation of a communication-heavy application on a Cray XE6 mainly to ``contention for
shared network resources from other jobs''~\cite{bhatele2013there}. Hoefler and Belli
found that only two of the 95 papers in their survey reported confidence intervals
around the mean, and they responded with twelve rules for measuring and reporting
performance~\cite{hoefler2015scientific}.

However, these rules leave choices to the experimenter, such as what to replicate, how
to time an exchange that ends at a barrier, and how to summarize per-rank times, and
those choices decide what the data can show. This document assesses the protocol of the
five-way campaign, whose results are reported separately~\cite{fivewayreport}, against
published guidance and against the project's research plan~\cite{researchplan}. For
each requirement we quote the guidance, describe what the campaign did with a citation to
the code that ran, and give a verdict of \emph{meets}, \emph{partly}, or \emph{does not}.
This document makes the following contributions:
\begin{itemize}
\item A description of the campaign's measurement protocol taken from the code that ran
  (\cref{sec:protocol});
\item An assessment of that protocol against published requirements, grouped by topic,
  with the consequence of each gap (\cref{sec:assessment});
\item A reusable checklist of those requirements, filled in for this campaign
  (\cref{sec:checklist});
\item A list of changes for the next campaign, each tied to the gap it closes
  (\cref{sec:changes}).
\end{itemize}
The assessment covers how the data was collected and summarized. Hoefler and Belli's
rules on presenting results (Rules 1, 3, 11, and 12) concern the report and are outside
its scope.

\section{The Protocol as Run}
\label{sec:protocol}

\paragraph{Units and kinds:} A \emph{launch} is one Slurm allocation. A \emph{run} is one
\texttt{srun} inside a launch, which starts a new MPI job for one combination of kind,
matrix, and arm. An \emph{iteration} is one halo exchange, or one CG iteration for the
CG kinds, inside a run. The campaign has seven kinds: three replays in the
irregular-benchmarking benchmark (the \pmrep{} of the exact matrix halo, the fine replay
of Vernier's per-rank capture, and the coarse replay of Vernier's aggregate histograms),
ABCG's SpMV exchange and ABCG's CG iteration with persistent requests, and the same two
ABCG kinds with one non-persistent call per exchange. The three arms are plain MPI
point-to-point (MPI), MPI Advance's neighbor alltoallv (STANDARD), and MPI Advance's
locality-aware neighbor alltoallv (LOCALITY). Each launch runs every combination of 7
kinds, 6 matrices, and 3 arms once, 126 runs.

\paragraph{Launches and order:} For each of 2, 4, 8, and 16 nodes, a pilot job and 40
launches were submitted, with each launch waiting for the launch four places before it,
so at most four launches of one node count ran at once (\path{scripts/submit.sh},
lines 10--12 and 48--52). Up to 16 launches ran at once across node counts, and the 160
launches spanned 23.7 hours. Each run is one \texttt{srun}
(\path{scripts/gatherdata.sbatch}, lines 170 and 176), and each launch runs its 126
runs in an order shuffled with a seed drawn from \texttt{/dev/urandom} and saved to the
launch's folder (lines 126--138), so the order of any launch can be reproduced. Launches
ran on the nodes Slurm assigned in the shared \texttt{pbatch} queue, with nodes dane261
and dane764 excluded.

\paragraph{Iteration counts:} The pilot ran each of the 42 patterns (kind and matrix) once
on STANDARD with $n=100$ and chose the iteration count $n$ that makes the $k+n$ iterations
total about 10\,s on STANDARD (\path{scripts/pilot.sbatch}, lines 148--160 for the
replays and 165--181 for ABCG). Every arm then uses that $n$. Every run performs $k=25$
iterations that are timed separately and excluded from the reported time per iteration.

\paragraph{Timing in the replay benchmark:} Before any pattern is measured, each rank
exchanges 64\,B and 128\,KiB messages with ranks $r\pm1$ on a duplicated communicator,
untimed (\path{src/benchmark.cpp}, lines 123--158). A barrier precedes the first timed
iteration (line 206). Each rank then times every iteration with
\texttt{std::chrono::high\_resolution\_clock} (line 34), from just before the exchange
until after a closing \texttt{MPI\_Barrier} (lines 213--219), and adds the elapsed time
to per-rank totals for the first $k$ and the next $n$ iterations. It also keeps the same
two totals stopped before the barrier, the time spent in the exchange alone (lines
215--227). Rank 0 receives the minimum, maximum, and sum over ranks of each total (lines
562--564) and prints the minimum, maximum, and mean; the analysis uses the maximum of the
barrier-inclusive total. Individual iteration times are not kept. In GCC's C++ library,
\texttt{high\_resolution\_clock} is an alias of \texttt{system\_clock}, a wall clock that
is not guaranteed to be monotonic.

\paragraph{Timing in ABCG:} ABCG's SpMV benchmark times each exchange barrier to barrier
with \texttt{MPI\_Wtime} (\path{benchmarks/benchmark_spmv.cpp}, lines 51--68). ABCG's
CG benchmark places one barrier before the first iteration and then records
\texttt{MPI\_Wtime} at the end of each iteration, so each iteration is timed from the end
of the previous one, with CG's own allreduces but no barrier in between
(\path{benchmarks/benchmark_cg_pilot.cpp}, lines 34--42 and 139--140); the
recomputation at the start of each new solve is charged to that solve's first iteration.
Neither ABCG benchmark keeps a per-rank time stopped before a synchronization point.

\paragraph{Correctness checks:} After timing, every replay run on STANDARD or LOCALITY
performs one extra, untimed exchange through both MPI Advance and plain MPI on a freshly
built halo and checks that every received value matches (\path{src/benchmark.cpp},
lines 613--620); all 5,760 such runs report \texttt{verify: PASS}. ABCG's SpMV benchmark
checks the buffer left by its last iteration against the expected columns
(\path{benchmarks/benchmark_spmv.cpp}, lines 73--77). ABCG's CG benchmark compares its
iteration count with ABCG's own solver, which is vacuous here: every completed solve
reached the 500-iteration cap, and in 360 runs no solve completed within $k+n$
iterations.

\paragraph{Provenance:} Each launch records its node list, the CPU and memory binding of
every rank, the MPI library, the Spack environment and its development checkouts, every
input's hash, and selected environment variables (\path{scripts/provenance.sh}); the
launch script saves the shuffle seed. Each replay run's output records its random seed,
MPI library, and compiler. ABCG's binaries are recorded only by hash, with their
compiler and flags stated in the campaign README. The replay benchmark's source is
recorded as ``not a git checkout'', so its commit is a hand-entered string, and the Spack
record's \texttt{locality\_aware} checkout (d10a789) is not the 752a19b the README says
the benchmark binary was built against. The provenance does not record the interconnect
(the fabric query on the compute nodes found no adapter it recognized), the processor
frequency settings, or the clock source, and Slurm's topology plugin on Dane is
\texttt{topology/flat}, so no switch placement is recorded.

\paragraph{Analysis:} The analysis treats the launch as the unit of replication. It
compares strategies by their medians over launches, with intervals from resampling whole
launches (10,000 resamples). It flags runs whose time exceeds five times their
combination's median, keeps them, and reports results with and without them. It
estimates position-in-launch effects and the variance shared by a launch with linear
mixed models.

\section{Assessment}
\label{sec:assessment}

\subsection{Replication}

Hunold and Carpen-Amarie found that ``it is insufficient for an MPI benchmark to collect
MPI run-time measurements only from a single call to mpirun''; ``10 calls to mpirun
seemed to be the minimum number of repetitions needed to obtain a good estimate of the
variance between mpiruns, while 30 repetitions were sufficient in all cases'', which
they offer as a rule of thumb~\cite[Section~4.2]{hunold2016reproducible}. Kalibera and
Jones state that ``At the very least, repetition must be done at the highest level that
has random variation to avoid bias''~\cite{kalibera2013rigorous}, and Georges et al.\
compute confidence intervals ``across the computed means from the different VM
invocations''~\cite{georges2007statistically}.

Every combination ran once in each of 40 launches per node count, each time as its own
\texttt{srun}, so each of its 40 measurements comes from a separate MPI job in a
separate allocation. Hunold and Carpen-Amarie's own repetitions were separate
\texttt{mpirun} calls inside one job (``All 30 calls to mpirun were executed as part of
the same compute job''), so the campaign replicates at a higher level than their
experiment and with more repetitions than their rule of thumb. The analysis computes one
statistic per run and its intervals across launches, as Georges et al.\ do across
invocations. Kalibera and Jones's highest level is the build, and the campaign used one
build of each program, so it meets their rule for the levels it varies but does not
replicate builds. \verdict{Meets for launches; partly for Kalibera and Jones's rule.}

\subsection{Randomization and Interleaving}

Hunold and Carpen-Amarie ``randomize the experiment by shuffling the order of
experiments within a call to mpirun''~\cite[Section~5.1]{hunold2016reproducible}, and
Hoefler and Belli write that ``If controlling a certain parameter is not possible then we
suggest randomization following standard textbook
procedures''~\cite{hoefler2015scientific}. Abedi and Brecht showed that single and
consecutive trials can ``falsely report that the performance of two identical systems
differ by 38\% using a confidence level of 95\%'', and that interleaving in a fixed order
fails when conditions change periodically, so ``trials must be randomized to avoid such
problems''~\cite{abedi2017conducting}. Their recommended design runs every alternative in
a fresh random order in each round. Maricq et al.\ recommend randomizing order ``to expose
effects that they might have on each other''~\cite{maricq2018taming}.

Each launch runs all 126 combinations in a freshly shuffled order, so with launches as
rounds the campaign follows Abedi and Brecht's design. The research plan asks for a
second design as well, rotating strategies iteration by iteration within one run for
claims about location~\cite{researchplan}. The campaign does not do this. Because every
strategy of a combination already runs in the same allocation, whole-allocation shifts
cancel in comparisons between strategies. However, the variation between runs within an
allocation does not cancel, and it is most of the variance, since a shift shared by the
whole launch accounts for only 9 to 23\% of the variance of log time~\cite{fivewayreport}.
\verdict{Meets for randomized order by round; does not interleave within a run.}

\subsection{Timing, Synchronization, and Per-Rank Summaries}

Hoefler and Belli's Rule~10 reads ``For parallel time measurements, report all
measurement, (optional) synchronization, and summarization techniques'';
\cref{sec:protocol} does that. Their Section~4.2.1 and the research plan make further
recommendations about how to time, which we take in turn.

On synchronization, Hoefler and Belli write that a barrier ``may be unreliable because
neither MPI nor OpenMP provides timing guarantees for their barrier calls. While a
barrier commonly synchronizes processes enough, we recommend checking the
implementation'', and they propose a window-based scheme. Hunold and Carpen-Amarie
conclude that ``For a fair comparison of MPI implementations, relying on MPI\_Barrier may
be completely sufficient if the same MPI\_Barrier algorithm is used by all of
them''~\cite[Section~4.3]{hunold2016reproducible}. In the replay and SpMV benchmarks every
arm's iterations close with the same Open MPI \texttt{MPI\_Barrier} on
\texttt{MPI\_COMM\_WORLD}, which meets their condition, and the barrier prevents the
``pipelining effects that occur when many operations are executed in a
row''~\cite{hoefler2008accurately}. The same 2008 paper notes that a barrier may introduce process skew and may use the same
network as the operation being measured, and the longer version of Hunold and Carpen-Amarie's study warns that
``processes can become significantly skewed when they leave
MPI\_Barrier''~\cite{hunold2015mpibenchmarking}. We did not measure barrier exit skew on
Dane. ABCG's CG benchmark has no barrier between iterations; its iterations are
separated only by CG's own allreduces.

On where the timer stops, the research plan's first measurement rule is to time each
rank ahead of every synchronization point, never around the iteration, because a timer
around a synchronized iteration measures the slowest rank at every rank~\cite{researchplan}.
The campaign's reported times include the closing barrier, so every rank's total includes
its wait for the slowest rank. Over the replay runs, the time outside the exchange
itself, which is the barrier wait plus any imbalance, is a median of 2 to 37\% of the
timed total, depending on node count and arm. The replay benchmark does keep per-rank
totals stopped before the barrier, but reports only their minimum, maximum, and mean;
the ABCG benchmarks keep none.

On summarizing across ranks, Hoefler and Belli write that a minimum or maximum across
processes ``seems sufficient'' in some work and ``We recommend avoiding such non-robust
measures'', proposing instead an ANOVA test of whether ranks' timings differ. The
campaign reports the maximum without such a test. Its barrier-inclusive totals agree
across ranks within 0.74\% in every run, but that agreement is produced by the barrier,
so it does not show that ranks behave alike.

On granularity, Hoefler and Belli write ``we recommend measuring single events to allow
the computation of confidence intervals and exact ranks''. The campaign keeps only totals
over $n$ iterations, so it has no per-iteration distribution and no within-run
percentiles, and it cannot tell whether a slow run took one long iteration or many
slightly slow ones. The 0.160\,s stalls reported for this campaign~\cite{fivewayreport}
could therefore be identified only from the quantization of per-run totals.

The replay benchmark's clock is \texttt{system\_clock}, which time synchronization can
adjust during a run. A clock adjustment would affect every arm alike, but only 12 of the
375 runs flagged as slow are LOCALITY runs, all of them non-persistent ABCG runs, so clock adjustments do not explain the
stalls. A monotonic clock would still be the safer choice. \verdict{Partly.}

\subsection{Warm-Up}

Hoefler and Belli note that communication systems establish state on demand, ``Thus, to
measure the expected time, the first measurement iteration should be excluded from the
average computation''~\cite[Section~4.1.2]{hoefler2015scientific}. Hunold and
Carpen-Amarie add that ``it is hard to estimate how many warm-up rounds are sufficient to
reach a `steady state'\,''~\cite[Section~3.3]{hunold2016reproducible}. The campaign
performs an untimed nearest-neighbor exchange before timing, excludes the first $k=25$
iterations from the reported time, and reports their total separately. We did not test
whether 25 iterations suffice for every combination. \verdict{Meets, with $k$
untested.}

\subsection{Number of Measurements and Run Length}

Hoefler and Belli recommend recomputing the confidence interval as measurements
accumulate and stopping once it is tight enough~\cite[Section~4.2.2]{hoefler2015scientific},
and Maricq et al.\ give a procedure for choosing repetition counts from observed
variability~\cite{maricq2018taming}. The campaign fixed 40 launches per node count in
advance and did not check interval widths before stopping. Some achieved intervals are
tight, such as the \pmrep's worst-case regret on CG, 19 to 26\%; others are not, such as
the fine replay's, 0 to 133\%, and the coarse replay's, 1,738 to
7,357\%~\cite{fivewayreport}. \verdict{Does not.}

Within a run, Hunold and Carpen-Amarie found that ``a sample size of 30 is not large
enough to obtain a normally distributed sample mean''~\cite[Section~4.1]{hunold2016reproducible},
which bears on how many iterations a per-run mean needs. The pilot set $n$ so a run takes
about 10\,s, but where the pilot itself stalled, $n$ came out small. At 16 nodes, the
pilot's STANDARD run of the fine replay of crankseg\_2 measured 142.8\,ms per iteration
and chose $n=57$. The fastest later MPI run of that pattern took 0.41\,ms per iteration,
so its timed loop lasts about 23\,ms rather than 10\,s, and one 0.160\,s stall adds
2.8\,ms to its time per iteration, about seven times the stall-free time; even the
fastest STANDARD run of that pattern (3.2\,ms) contains one stall. Across the campaign,
the pilot's time per iteration exceeded the fastest later STANDARD run's by more than 2
times for 47 of the 168 combinations of node count and pattern, and by more than 10 times
for 19 of them. \verdict{Partly.}

\subsection{Statistics}

Hoefler and Belli's Rule~5 asks for confidence intervals for nondeterministic data, Rule~6
says ``Do not assume normality of collected data (e.g., based on the number of samples)
without diagnostic checking'', and Rule~7 asks that comparisons be made ``in a
statistically sound way''~\cite{hoefler2015scientific}. They also note that
``Nonparametric metrics, such as the median or other percentiles, do not assume a specific
distribution and are most robust'', and add that such measures still require independent
and identically distributed measurements. Maricq et al.\ advise ``Use nonparametric
confidence intervals to avoid assumptions of normality''~\cite{maricq2018taming}, and
Agarwal et al.\ recommend bootstrap interval estimates when only a few runs are
available~\cite{agarwal2021deep}.

Every comparison between strategies uses medians over launches and intervals from
resampling whole launches, with no assumed distribution, and ratios are compared with a
bootstrap equivalence test at a $\pm10\%$ margin fixed before computing. However, the
position-in-launch effect and the variance shares come from linear mixed models with
normally distributed errors on log time, whose tails the stalls inflate, and we ran no
normality diagnostic for them. Rule~4 reads ``Avoid summarizing ratios; summarize the
costs or rates that the ratios base on instead. Only if these are not available use the
geometric mean for summarizing ratios.'' The campaign's selection-gain statistic is a
geometric mean of ratios across matrices, chosen because CG's times per iteration
differ across matrices and arms by 11 to 147 times at a single node count; the costs
were available, so this departs from Rule~4. Rule~8
reads ``Carefully investigate if measures of central tendency such as mean or median are
useful to report. Some problems, such as worst-case latency, may require other
percentiles.'' The analysis reports the 10th percentile and robust spread over launches,
the former as an exploratory check added after the stalls were found, but it has no
within-run percentiles. \verdict{Partly.}

\subsection{Outliers and Anomalies}

Hoefler and Belli write ``We suggest to avoid removal of outliers and instead use robust
measures such as percentiles'', and, where removal is necessary, ``one should report the
number of removed outliers for each experiment''~\cite{hoefler2015scientific}. The
analysis keeps every run, uses medians, flags runs over five times their combination's
median, reports the number flagged (266 of 9,600 MPI and STANDARD runs of the five kinds
the report studies), and shows each result with and without them. The five-times flag
misses most stall-affected runs where stalls inflate the median itself, so results
without flagged runs are only a partial check. \verdict{Meets.}

\subsection{Interference, Stationarity, and Independence}

Bhatele et al.\ found ``the primary reason for higher variability is contention for
shared network resources from other jobs''~\cite{bhatele2013there}. Skinner and Kramer
list contention for shared resources and system transitions among the causes of
variability~\cite{skinner2005understanding}. Duplyakin et al.\ argue that measurements on
clouds, HPC facilities, and datacenters are ``almost guaranteed to exhibit performance
regimes that evolve over time'', and found such changes in 2.75 years of CloudLab
benchmark data~\cite{duplyakin2020datacenter}. Kalibera and Jones warn that ``If
measurements are not i.i.d., the variance and confidence interval estimates will be
biased''~\cite{kalibera2013rigorous}.

The campaign ran in a shared queue and treats interference from other jobs as part of the
variation between launches. Neither launch start time (Spearman $|\rho| \le 0.12$,
$p \ge 0.45$ at every node count) nor membership of any one node (smallest
Bonferroni-corrected $p = 0.73$ over 124 nodes) showed an association with which launches
stall, and the position effect within a launch is significant only at 8
nodes~\cite{fivewayreport}, by the mixed model noted above. The campaign did not test
whether its own concurrent launches, up to 16 at once, interfere with each other, and its
23.7-hour span cannot show changes over weeks, for which the research plan asks for
allocations spread across weeks~\cite{researchplan}. \verdict{Partly.}

\subsection{Documentation and Provenance}

Hoefler and Belli's Rule~9 reads ``Document all varying factors and their levels as well
as the complete experimental setup (e.g., software, hardware, techniques) to facilitate
reproducibility and provide interpretability'', and their Section~4.1.2 asks for network
details. Hunold and Carpen-Amarie write that ``If an MPI benchmarking experiment does not
clearly state the compiler and the compilation flags used, the results will not be
comparable or might not even be trustworthy''~\cite[Section~4.5]{hunold2016reproducible},
and they show that pinning and the processor frequency setting change MPI timings
(Sections 4.4 and 4.6). Dodge et al.\ include a description of the computing
infrastructure in their reporting checklist~\cite{dodge2019show}. Rule~2 asks
experimenters to ``Specify the reason for only reporting subsets of standard benchmarks
or applications or not using all system resources''.

For the replay benchmark, the campaign records per launch the binary's hash, compiler,
flags, and MPI library; for ABCG it records only binary hashes, with compiler and flags
stated in the campaign README. The benchmark's commit is entered by hand, and the Spack
record's \texttt{locality\_aware} checkout does not match the commit the binary was built
against, so the provenance does not by itself identify the MPI Advance code that ran. The
campaign records per-rank binding but not the interconnect, the processor frequency
settings, or the clock source, and its documents do not state why these six SuiteSparse
matrices were chosen. \verdict{Partly; does not for Rule~2.}

\subsection{Correctness of What Is Measured}

None of the papers above makes this a rule, but a strategy that delivers wrong data is
not a valid point of comparison. The replay benchmark's check covers one untimed exchange
per run after timing, against plain MPI, and all 5,760 checked runs passed. ABCG's SpMV
check covers the last iteration's buffer. ABCG's CG check is vacuous, so no CG run has an
effective correctness check. \verdict{Partly.}

\section{Checklist}
\label{sec:checklist}

\Cref{tab:checklist} lists each requirement with its source, what this campaign did, and
the verdict. A future campaign can keep the first two columns and fill in the last two.

{\small
\setlength{\tabcolsep}{4pt}
\begin{longtable}{@{}>{\raggedright\arraybackslash}p{0.05\linewidth}>{\raggedright\arraybackslash}p{0.33\linewidth}>{\raggedright\arraybackslash}p{0.44\linewidth}>{\raggedright\arraybackslash}p{0.10\linewidth}@{}}
\caption{Requirements and this campaign's protocol. H\&B: Hoefler and
Belli~\cite{hoefler2015scientific}; H\&C: Hunold and
Carpen-Amarie~\cite{hunold2016reproducible}; Plan: the project's research
plan~\cite{researchplan}.}
\label{tab:checklist}\\
\toprule
\# & Requirement (source) & This campaign & Verdict \\
\midrule
\endfirsthead
\toprule
\# & Requirement (source) & This campaign & Verdict \\
\midrule
\endhead
1 & Repeat across separate \texttt{mpirun} calls, 10 to 30 as a rule of thumb (H\&C 4.2) & 40 launches per node count, each run its own \texttt{srun} & Meets \\
2 & Repeat at the highest level with random variation (Kalibera and Jones) & Launches replicated; builds not replicated & Partly \\
3 & Intervals across invocations, not iterations (Georges et al.) & One statistic per run, intervals across launches & Meets \\
4 & Randomize order within a launch (H\&C 5.1; H\&B 4.1.1) & Order shuffled per launch with a saved seed & Meets \\
5 & Randomized interleaving in rounds (Abedi and Brecht) & Each launch a round with a fresh random order & Meets \\
6 & Interleave strategies per iteration for location claims (Plan 5.4) & Not done & Does not \\
7 & Report timing, synchronization, and summarization (H\&B Rule~10) & Stated in \cref{sec:protocol} & Meets \\
8 & Measure single events (H\&B 4.2.1) & Totals over $n$ iterations only & Does not \\
9 & Time each rank ahead of synchronization (Plan 5.4) & Reported times include the barrier; pre-barrier totals kept only as min, max, and mean for the replays; none for ABCG & Partly \\
10 & \texttt{MPI\_Barrier} suffices when every alternative uses the same barrier (H\&C 4.3) & Same barrier for every arm; exit skew not measured & Partly \\
11 & Test whether ranks differ before summarizing (H\&B 4.2.1) & Maximum used without a test & Does not \\
12 & Avoid pipelining between consecutive operations (Hoefler et al.\ 2008) & Barrier after every replay and SpMV iteration; CG iterations separated by its allreduces & Meets \\
13 & Exclude warm-up (H\&B 4.1.2; H\&C 3.3) & Untimed warm-up plus 25 excluded iterations; adequacy untested & Meets \\
14 & Stop when intervals are tight enough (H\&B 4.2.2; Maricq et al.) & 40 launches fixed in advance & Does not \\
15 & Enough iterations per measurement (H\&C 4.1) & $n$ from pilots that ran more than 2 times slower than the fastest later run in 47 of 168 cases & Partly \\
16 & Confidence intervals for nondeterministic data (H\&B Rule~5) & Bootstrap intervals over launches & Meets \\
17 & No normality without checking (H\&B Rule~6; Maricq et al.) & Nonparametric for comparisons; normal mixed models unchecked & Partly \\
18 & Statistically sound comparison (H\&B Rule~7) & Bootstrap intervals and equivalence test & Meets \\
19 & Summarize costs, not ratios (H\&B Rule~4) & Selection gain is a geometric mean of ratios & Partly \\
20 & Consider percentiles beyond the median (H\&B Rule~8) & 10th percentile and spread over launches; none within runs & Partly \\
21 & Keep outliers; report how many are removed (H\&B 3.1.3) & All kept, flagged, shown with and without; flag is a partial check & Meets \\
22 & Account for interference and nonstationarity (Bhatele et al.; Duplyakin et al.) & Shared queue; 23.7 hours; concurrent launches untested & Partly \\
23 & Check independence (Kalibera and Jones; Maricq et al.) & Start time, node, and position tested & Partly \\
24 & Document the complete setup (H\&B Rule~9) & Most recorded; interconnect, frequency, clock source missing; benchmark commit hand-entered & Partly \\
25 & State compiler and flags (H\&C 4.5) & Per launch for the replays; README only for ABCG & Partly \\
26 & Record pinning and frequency settings (H\&C 4.4, 4.6) & Binding recorded; frequency not & Partly \\
27 & Justify subsets of benchmarks (H\&B Rule~2) & Reason for the six matrices not stated & Does not \\
28 & Verify delivered data (not a published rule) & One untimed exchange per replay run; last SpMV buffer; CG vacuous & Partly \\
\bottomrule
\end{longtable}
}

Of the 28 requirements, the campaign meets 10, partly meets 13, and does not meet 5.

\section{Changes for the Next Campaign}
\label{sec:changes}

There are several changes that address the gaps above, in order of how much they would
change what the data can show. First, the benchmarks should record per-iteration,
per-rank timestamps in memory, stopped before each synchronization point, and write them
at the end of each run (items 8, 9, 11, and 20). This would show directly whether a stall
is one 0.160\,s iteration and would allow a test of whether ranks differ. Second, claims
about location should come from runs that interleave strategies iteration by iteration
(item 6), with the whole-run design kept for spread between allocations. Third, the pilot
should choose $n$ from a stall-robust estimate, such as the median of several pilot runs,
with a floor on $n$ (item 15). Fourth, the campaign should check interval widths as
launches accumulate and stop on a stated target (item 14). Fifth, launches should be
spread over weeks and across at least one system change, and a controlled test should
check whether concurrent launches of the campaign interfere (items 22 and 23). Sixth,
provenance should come from git checkouts of every program, a Spack record that matches
the builds used, the compiler and flags of every binary, and the interconnect, frequency
setting, \texttt{MPI\_Wtick}, and clock source, and the benchmark should use
\texttt{std::chrono::steady\_clock} (items 24 to 26). Seventh, the campaign documents
should state why each matrix was chosen (item 27), the analysis should summarize costs
rather than ratios where costs exist (item 19), and the mixed models should be checked
against a normality diagnostic or replaced by a resampling method (item 17). Eighth, the
CG benchmark should use a tolerance its solves can reach so its correctness check is not
vacuous, and the replay check should cover a timed exchange (item 28). Finally, a short
experiment should measure barrier exit skew on Dane (item 10).

\section{Conclusion}

In this document we assessed the five-way campaign's data-collection protocol against 28
requirements drawn from Hoefler and Belli, Hunold and Carpen-Amarie, related HPC,
systems, and machine learning work, and the project's research plan. The campaign
replicated every combination in 40 separate allocations per node count, randomized run
order within each allocation with recorded seeds, compared strategies with medians and
intervals resampled over allocations, and kept and reported every outlier. It met 10
requirements, partly met 13, and did not meet 5. It did not time single iterations, did
not interleave strategies within a run, did not test whether ranks differ before taking
their maximum, fixed its number of allocations in advance, and did not state why its six
matrices were chosen. It also included a barrier inside every timed interval of the
replay and SpMV benchmarks, chose iteration counts from pilots that ran more than 2 times
slower than the fastest later run in 47 of 168 cases, spanned 23.7 hours rather than weeks, and left gaps in its provenance and
correctness checks. The checklist in \cref{sec:checklist} can be refilled for the next
campaign.

\bibliographystyle{plain}
\bibliography{references}

\end{document}
